\section{Harness Engineering 与 Context Engineering}

\subsection{Harness Engineering 的定义}

Harness engineering 这个术语由 Viv 提出，其核心思想是：\textbf{每当你发现 agent 犯了一个错误，就花时间工程化一个解决方案，使得 agent 永远不会再犯同样的错误。}（Mitchell Hashimoto 的总结）

\begin{importantbox}{Harness Engineering 的六个核心问题}
\begin{enumerate}[leftmargin=1.5em]
    \item 如何赋予 coding agent 新的能力？
    \item 如何教会它训练数据中没有的代码库知识？
    \item 如何在 system message 中的 ``CRITICAL: always do XYZ'' 之外增加确定性？
    \item 如何让 agent 的行为适配特定代码库？
    \item 如何在``魔法 prompt''之外提高任务成功率？
    \item 如何防止 context window 过快膨胀或被低质量上下文填满？
\end{enumerate}
\end{importantbox}

\subsection{Context Engineering 的超集关系}

HumanLayer 的联合创始人 Dex 在 12-factor agents 中提出了 context engineering 的概念。Context engineering 是 prompt engineering 的超集，涵盖了系统性提高 AI agent 可靠性的各种技术。

Harness engineering 则是 context engineering 的子集，\textbf{主要通过 harness 配置点来精心管理 coding agent 的 context window}。

\begin{knowledgebox}{Viv 的四个定制杠杆}
Viv 关于 harness engineering 的文章将定制杠杆分为四类：
\begin{enumerate}[leftmargin=1.5em]
    \item \textbf{System prompt}：系统级指令
    \item \textbf{Tools / MCPs}：工具与能力扩展
    \item \textbf{Context}：上下文管理
    \item \textbf{Sub-agents}：子代理分工
\end{enumerate}
HumanLayer 额外强调了两个 Viv 未重点讨论的杠杆：\textbf{hooks}（自动化集成与确定性控制流）和 \textbf{skills}（渐进式知识披露）。
\end{knowledgebox}

\subsection{Post-Training 与 Harness 的耦合}

前沿 coding 模型会在其对应 harness 上进行 post-training（例如 Claude 在 Claude Code 上训练，GPT-5 Codex 在 Codex harness 上训练）。这导致了有趣的现象：

\begin{warningbox}{Post-Training 的双刃剑}
模型对其训练 harness 的过度拟合可能适得其反。Terminal Bench 2.0 的数据显示：Opus 4.6 在 Claude Code 中排名第 33，但放到一个 post-training 时未见过的 harness 中，排名升至第 5（上下浮动约 4 名）。这说明 \textbf{默认 harness 不一定是最优 harness}，定制化配置可能释放模型被默认 harness 压制的潜力。
\end{warningbox}

一个典型案例是 OpenCode 项目。Codex 模型与 Codex harness 的 \texttt{apply\_patch} 工具紧密耦合，以至于 OpenCode（一个开源的 Claude Code 替代方案）不得不专门为 GPT/Codex 模型添加 \texttt{apply\_patch} 工具来模拟 Codex harness，而 Claude 和其他模型则继续使用常规的 edit 和 write 工具。

\subsection{本章小结}

Harness engineering 是 context engineering 的子集，通过配置 harness 的各个组件来管理 agent 的 context window。核心理念是：遇到错误就工程化一个永久性修复。即使模型经过了 post-training，默认 harness 也并非不可改进。


